\twocolumn
\setlength{\emergencystretch}{3em}

\section{The one law, and the four working
constraints}\label{the-one-law-and-the-four-working-constraints}

Most physical theories begin with many laws: one for electricity, one
for gravity, one for heat, plus a list of particles to put on the stage.
The Space-Conversion Programme (SCP) begins with one:

\begin{quote}
\textbf{Energy is never destroyed. It only changes mode. Space is one of
the modes.}
\end{quote}

Everything else in this paper is either a \emph{consequence} of that
sentence, a \emph{definition} needed to state it precisely, or a
\emph{measurement} of how well the consequences hold up in simulation.
Four working constraints keep the construction honest:

\begin{enumerate}
\def\labelenumi{\arabic{enumi}.}
\tightlist
\item
  \textbf{Energy is never destroyed} --- the only law. In the simulation
  this is not a hope but a checked number: the total energy of a closed
  run drifts only at the level of computer round-off (parts in
  \(10^{13}\)--\(10^{15}\)), because every operation in the update
  conserves energy \emph{by construction} (§3.2, §8).
\item
  \textbf{No background.} Nothing in the theory is allowed to persist
  forever and merely have its values updated. There is no fixed grid of
  points, no eternal coordinate system. The cells of §2 are created and
  destroyed; space is a \emph{participant}, not a stage.
\item
  \textbf{No imported species.} We do not add an ``electron field'' or a
  ``photon'' by hand. Whatever objects exist must \emph{emerge} as
  configurations of the substrate.
\item
  \textbf{Emergent from the fabric alone.} Familiar laws (Maxwell,
  \(E=mc^2\), the \(T^4\) radiation law, Newtonian gravity, quantum
  interference) are targets to be \emph{derived and measured}, never
  assumptions.
\end{enumerate}

\textbf{How claims are graded.} Throughout, statements carry one of four
tags. \textbf{{[}derived{]}}: follows mathematically from the substrate
rules. \textbf{{[}measured{]}}: confirmed in simulation, with the number
quoted. \textbf{{[}adopted{]}}: a rule or constant chosen and then held
fixed under a regression battery (the ``ratchet'': every change to the
substrate must re-pass every previously passing experiment --- 87 gated
experiments at the current law table). \textbf{{[}open{]}}: honestly
unresolved. A theory paper that never says ``open'' is hiding something;
§10 is a list of what we are hiding from no one.

\textbf{Two ways to read.} The mathematical spine runs §3 → §4 → §5;
every equation in it is optional in the sense of §2.3 (the words carry
the content). For the story without the equations, read §§1--2, §4.1,
§4.4, §6, §7, the plain-language theorem at the start of §8.2, and §9
--- whose correspondence table (Table \ref{tab:corr}) can equally be
read \emph{now}, as a map of where this paper is going.

\section{The substrate: cells, modes, and the
door}\label{the-substrate-cells-modes-and-the-door}

\subsection{Cells and the roster}\label{cells-and-the-roster}

The world is a finite collection of \textbf{cells}. A cell is not a
location in space --- it is a \emph{piece} of space: a parcel of the
space mode of energy, with a size set by how much space-energy it holds.
Cells that are close enough touch; touching cells are joined by
\textbf{links}, along which energy can be exchanged. The set of live
cells and links --- the \textbf{roster} --- is itself dynamical: cells
are born when energy condenses, die when they drain, and the linkage
rearranges as cells move and resize. This is what ``no background''
means operationally: \emph{there is no list of places; there is only the
roster, and the roster is made of the same stuff it describes.}

\subsection{The three sectors and the
register}\label{the-three-sectors-and-the-register}

\begin{figure*}[t]
\centering
\includegraphics[width=0.92\textwidth]{fig_substrate.pdf}
\caption{\textbf{The substrate.} (a) One cell carries three energy modes
--- field $\psi_e$, matter $\psi_m$, space $E_s$ --- plus the register
$\chi$, which holds energy \emph{in the act of converting} between field
and matter. The field--register connection is continuous; the
register--matter connection fires only in discrete clicks of whole atoms.
(b) The full graph: spatial field links, spatial matter links, and the
per-cell mode (``fiber'') links. One rotation primitive operates on every
edge type.}
\label{fig:substrate}
\end{figure*}

Each cell holds energy in three \textbf{modes} (sectors), plus one
bookkeeping carrier (Fig. \ref{fig:substrate}):

\begin{itemize}
\tightlist
\item
  \textbf{The field}, written \(\psi_e\). This is the light-like sector:
  it is fast, it carries waves, and in the vacuum it travels at the
  substrate's top speed. \(\psi_e\) is a \emph{complex amplitude} ---
  explained below.
\item
  \textbf{Matter}, written \(\psi_m\), with energy \(E_m\). This is the
  slow, storable sector: the stuff lumps, objects, and (eventually)
  particles are made of.
\item
  \textbf{Space}, written \(E_s\). A plain (non-oscillating) energy
  store. The amount of \(E_s\) a cell holds sets its size. Matter is
  \emph{converted space}: making matter somewhere means space-energy was
  consumed there, and that deficit is the seed of gravity (§6).
\item
  \textbf{The register}, written \(\chi\): a small complex amplitude
  holding energy \emph{in transit between modes} --- energy that has
  left the field but has not yet become matter (or the reverse). It is
  the exact cousin of the ``flight'' parcels that carry energy mid-link
  between two cells (Fig. \ref{fig:transit}): flight is energy in
  space-transit; the register is energy in \emph{mode}-transit. It is
  bookkeeping for an incomplete conversion, not a fourth kind of
  substance --- a claim we make precise, and prove necessary, in §8.
\end{itemize}

\begin{figure}[t]
\centering
\includegraphics[width=\linewidth]{fig_transit.pdf}
\caption{\textbf{Energy in transit, twice.} (a) A flight parcel: energy
mid-hop on a link between two cells --- it has left cell A and not yet
arrived at cell B, and the ledger books it separately. (b) The register
$\chi$: the same idea in the \emph{mode} direction --- energy that has
left the field and not yet settled as matter, inside one cell. Both are
counted in the conservation sum; neither is a new kind of substance.}
\label{fig:transit}
\end{figure}

\textbf{What a complex amplitude is.} A complex amplitude is just an
arrow in a plane: it has a \emph{length} and a \emph{direction} (a
phase, i.e.~a clock reading). We write it as two ordinary numbers ---
\(\psi_e = f_{a1} + i f_{a2}\) means ``an arrow whose two perpendicular
components are \(f_{a1}\) and \(f_{a2}\).'' The \textbf{energy} of the
sector is the squared length of the arrow:
\[E_e = |\psi_e|^2 = f_{a1}^2 + f_{a2}^2 .\] The direction of the arrow
is the sector's internal clock. Nothing quantum is being smuggled in
here: an arrow with a length and a clock is exactly how one describes
any oscillation (a pendulum has an amplitude and a phase). What
\emph{will} be derived --- not assumed --- is that adding arrows before
squaring produces interference (§7). (Every symbol in this paper is
collected in Appendix A; dip into it freely.)

\subsection{The arrow, the quarter-turn, and
energy}\label{the-arrow-the-quarter-turn-and-energy}

Because the rest of the paper leans on this one picture, we spell it out
once, slowly (Fig. \ref{fig:arrow}). An amplitude is an arrow. Its
\emph{length} says how big the oscillation is; its \emph{angle} is a
clock hand saying where in its cycle the oscillation is. Two operations
cover everything the substrate ever does to arrows:

\begin{itemize}
\tightlist
\item
  \textbf{Turning.} Multiplying by the imaginary unit \(i\) is nothing
  deeper than \emph{turning the arrow a quarter-turn} (\(90^\circ\))
  counter-clockwise. That is the only sense in which \(i\) appears in
  this paper: it is a rotation instruction, not a mysterious number. An
  equation like
  ``\(\text{rate of change of } \psi = -i\,\Omega\,\psi\)'' says, in
  words: \emph{the arrow's tip moves at right angles to the arrow} ---
  which is exactly what ``rotating steadily at rate \(\Omega\)'' means.
\item
  \textbf{Squaring the length.} The energy carried by the arrow is its
  length squared, \(E = f_{a1}^2 + f_{a2}^2\) --- always positive, zero
  only for a zero arrow. Two arrows can point opposite ways and
  \emph{add} to a short arrow (that is interference, §7), but energy is
  only ever computed from the arrow a cell actually has --- never by
  cancelling one cell's energy against another's.
\end{itemize}

Readers comfortable with complex numbers can read every equation below
as written; readers who are not can, everywhere, substitute the words
``turn a quarter-turn'' for \(i\) and ``rotate steadily'' for
\(e^{-i\Omega t}\), and no content is lost.

\begin{figure*}[t]
\centering
\includegraphics[width=0.92\textwidth]{fig_arrow.pdf}
\caption{\textbf{The one picture behind all the math.} (a) An amplitude
is an arrow with a length and a clock angle. (b) The imaginary unit $i$
is a quarter-turn: $i\psi$ is $\psi$ rotated $90^\circ$. (c) Energy is
the squared length of the arrow --- the sum of the squares of its two
perpendicular components, and never negative.}
\label{fig:arrow}
\end{figure*}

\subsection{Pitches, load, and the flattening
law}\label{pitches-load-and-the-flattening-law}

Every oscillating sector has a natural frequency, its \textbf{pitch}:
\(w_1\) for the field and \(w_2\) for matter, with \(w_2 > w_1\)
(matter's clock runs faster than the field's --- in the adopted law
table, \(w_1 = 1.65\), \(w_2 = 2.9\) in substrate units). Pitch is not
fixed: a loaded cell runs slow. Define the \textbf{load}
\[x = \frac{E_m + \text{(energy in transit)}}{\mathrm{cap}},\] the
fraction of the cell's carrying capacity \(\mathrm{cap}\) that is in
use. Then the effective pitches are \textbf{{[}adopted,
battery-gated{]}} \[w_{1e} = \frac{w_1}{1 + q\,x}, \qquad
  w_{2e} = \frac{w_2}{1 + q\,x},\] with a single constant \(q\)
(\(=1.2\) in the law table). \emph{Load flattens pitch.} This one rule
--- a full cell ticks slower --- will reappear as refraction (§5.2), as
the work function of the conversion door (§5.1), and as a click-rate
throttle (§4.3).

\subsection{The door and the atoms}\label{the-door-and-the-atoms}

Energy moves \emph{within} a mode continuously. Energy moves
\emph{between} modes only through the \textbf{door}, and the door is
quantized: every conversion fires in whole \textbf{action atoms}
\[\boxed{\;\varepsilon = \frac{A_0\,\omega}{2\pi}\;}\] --- the energy
that carries exactly \(A_0\) of \emph{action} (energy \(\times\) time)
in one cycle of the source's pitch \(\omega\). Why that form: one full
turn of a clock ticking at pitch \(\omega\) takes time \(2\pi/\omega\),
so a packet that carries action \(A_0\) in one turn carries energy
\(A_0 \div (2\pi/\omega) = A_0\omega/2\pi\) --- the \(2\pi\) is just
``per full turn.'' The constant \(A_0\) (\(=1.15\) in the law table) is
the substrate's version of Planck's constant: the size of the smallest
tradable packet of oscillation. That conversion is atom-grained while
transport is smooth is the deepest structural fact of the substrate ---
``amplitudes within a mode, atoms at mode boundaries'' --- and §8 proves
the two cannot be merged. Linearity of \(\varepsilon\) in \(A_0\) is
\textbf{{[}measured{]}} to about one part in \(10^{8}\).

\subsection{The conservation ledger}\label{the-conservation-ledger}

The books close over everything:
\[\frac{d}{dt}\;\sum_{\text{cells}}\Big(E_m + E_e + E_s + |\chi|^2\Big)
  \;+\; \frac{d}{dt}\,\big(\text{flight}\big) \;=\; 0 .\]
\textbf{{[}derived + measured{]}} Derived, because (as §3--§4 show)
every continuous update step is a rotation --- and rotations do not
change lengths --- while every discrete click is an explicit ledger
move; measured, because closed runs hold this sum to
\(\sim 10^{-13}\)--\(10^{-15}\) relative drift over thousands of time
units, and a single bookkeeping bug (an uncapped clamp in one kernel's
outtake branch) announced itself as a drift eight orders of magnitude
above that floor and was found \emph{by this meter}.

\section{The translation primitive: how anything
moves}\label{the-translation-primitive-how-anything-moves}

\subsection{Two half-oscillations that feed each
other}\label{two-half-oscillations-that-feed-each-other}

How does a wave move? In the substrate, the field update per time step
\(dt\) consists of exactly two operations. \textbf{{[}derived from the
kernel; this section is the measured code, written as mathematics{]}}

\textbf{(a) Precession.} Each cell's field arrow rotates in place at its
effective pitch: \(\psi_e \to e^{-i\,\Omega_1 dt}\,\psi_e\) with
\(\Omega_1 = w_{1e}\). Component-wise, the arrow \((f_{a1}, f_{a2})\)
turns by angle \(\Omega_1 dt\).

\textbf{(b) Exchange.} Each link mixes the arrows of its two end cells
by a small energy-preserving rotation (a \emph{Givens rotation} --- the
plane rotation familiar from geometry, applied to the pair of
amplitudes): \[\psi_i \to \cos\tau\,\psi_i - i\sin\tau\,\psi_j, \qquad
  \psi_j \to \cos\tau\,\psi_j - i\sin\tau\,\psi_i,\] with a small angle
\(\tau\) per step. Because it is a rotation, \(|\psi_i|^2 + |\psi_j|^2\)
--- the two cells' combined energy --- is \emph{exactly} unchanged.
Conservation is not imposed; it is geometry. (The \(-i\) in the mix rule
says the neighbour's arrow arrives \emph{turned a quarter-turn} --- Fig.
\ref{fig:arrow}b --- the same \(90^\circ\) relationship the arrow's own
two components keep between themselves. Those two perpendicular
components, \(f_{a1}\) and \(f_{a2}\), are called the
\textbf{quadratures}.)

Now write the same two operations in terms of the arrow's two
components. With neighbours indexed by \(j\) and exchange strengths
\(K_{ij}\):
\[\frac{d f_{a1}(i)}{dt} = +\Omega_1 f_{a2}(i) + \sum_j K_{ij} f_{a2}(j),\]
\[\frac{d f_{a2}(i)}{dt} = -\Omega_1 f_{a1}(i) - \sum_j K_{ij} f_{a1}(j).\]
Look at the \emph{shape} of this pair, because it is the whole point:

\begin{itemize}
\tightlist
\item
  each equation is \textbf{first-order} (no accelerations anywhere);
\item
  each is \textbf{incomplete} --- \(f_{a1}\)'s rate of change contains
  no \(f_{a1}\) at all: it is sourced \emph{entirely} by the other
  component, on-site and on the neighbours; and symmetrically for
  \(f_{a2}\);
\item
  the two are \textbf{mutually inducing}: freeze either component and
  the other goes inert.
\end{itemize}

Propagation is the \emph{interleaving} of two incomplete conversions
(Fig. \ref{fig:maxwell}). Neither half moves on its own: precession
alone is oscillation in place; exchange alone spreads a lump
symmetrically with no preferred direction. A travelling wave is a phase
\emph{gradient} (the arrows' clocks reading slightly differently from
cell to cell) advanced by precession while exchange converts that
gradient into displacement. Figure \ref{fig:pulsehop} shows it
happening: these are not sketches but the actual update run on a 40-cell
chain --- precession, then the pairwise exchange sweep, repeated --- and
the pulse walks.

\begin{figure}[t]
\centering
\includegraphics[width=\linewidth]{fig_pulsehop.pdf}
\caption{\textbf{Watch the pulse move.} Four snapshots of the actual
precession-plus-exchange update on a 40-cell chain (energy per cell,
successive times). No cell ever ``sends'' anything: each step only
rotates neighbouring pairs of arrows, yet the energy packet advances at
the band speed. Total energy is identical in every frame.}
\label{fig:pulsehop}
\end{figure}

\begin{figure*}[t]
\centering
\includegraphics[width=0.92\textwidth]{fig_maxwell.pdf}
\caption{\textbf{The translation primitive.} (a) In one cell the two
quadratures $f_{a1}, f_{a2}$ form a rotating arrow; each component's rate
of change is sourced only by the other. (b) Across many cells the same
pairing propagates: the two components run a quarter-cycle apart and hand
the waveform forward at the band speed $c_{\mathrm{sub}}$. This is the
Maxwell shape --- two first-order, mutually-inducing, incomplete
conversions --- realized on the roster.}
\label{fig:maxwell}
\end{figure*}

\subsection{The Maxwell
correspondence}\label{the-maxwell-correspondence}

\begin{quote}
\textbf{The familiar law.} In everyday electromagnetism, light is two
fields --- electric \(\mathbf{E}\) and magnetic \(\mathbf{B}\) --- and
the essential content of Maxwell's equations in empty space is a loop:
\emph{the electric field's rate of change is caused by the magnetic
field, and the magnetic field's rate of change is caused by the electric
field} (with one sign flipped). Neither field sustains itself; each
exists by inducing the other, and the loop propagates at a fixed speed
\(c\). For readers who have met the symbols, the pair is
\[\frac{\partial \mathbf{E}}{\partial t} = c^2\,\nabla\times\mathbf{B},
  \qquad
  \frac{\partial \mathbf{B}}{\partial t} = -\,\nabla\times\mathbf{E};\]
the words above are the content, and the rest of this section needs only
the words.
\end{quote}

Compare with the substrate pair above. The correspondence is structural
and exact in the properties that matter: two first-order equations; no
self-terms; mutual induction with one sign flip; elimination of either
component gives a second-order wave equation \(d^2\psi/dt^2 = -H^2\psi\)
on the roster. In this sense \textbf{the field sector does not
\emph{satisfy} a Maxwell-like law; it \emph{is} the Maxwell mechanism},
with \(f_{a1}\) and \(f_{a2}\) playing the roles of \(\mathbf{E}\) and
\(\mathbf{B}\).

One honest caveat \textbf{{[}stated, not hidden{]}}: the textbook
\(\nabla\times\) acts on \emph{vector} fields in three dimensions; the
roster field is one complex amplitude per cell, so the neighbour
coupling is the graph operator \(K\), not a literal curl. What carries
over is the shape --- and the shape is what forces wave behaviour.

\emph{(Skippable remark, for readers who have met quantum mechanics;
everyone else may jump to §3.3 losing nothing.)} The same pair, written
for the complex arrow, is a discrete Schrödinger equation
\[i\,\frac{d\psi}{dt} = H\psi,\] where the matrix
\(H = \mathrm{diag}(\Omega_1(i)) + K\) is symmetric. That the
\emph{classical field update} and the \emph{Schrödinger form} are the
same object on the substrate is not an analogy --- they are literally
the same two lines of arithmetic.

\subsection{The speed of light,
derived}\label{the-speed-of-light-derived}

For a wave \(\psi \propto e^{i(kx - \omega t)}\) on a regular chain of
cells with spacing \(d\) and mean exchange rate \(\bar K\), the update
above gives the \textbf{dispersion relation} \textbf{{[}derived{]}}
\[\omega(k) = \Omega_1 + \epsilon(k), \qquad
  \epsilon(k) = 2\bar K\,(1 - \cos kd),\] and a packet travels at the
\emph{group velocity} \(v_g = d\omega/dk\), whose maximum over the band
is the substrate speed of light:
\[\boxed{\;c_{\mathrm{sub}} \;\approx\; 2\,\bar K\, d\;}\] fixed
entirely by the law constants and the roster geometry --- no speed is
imported. Every wave in the vacuum, regardless of wavelength or
intensity, is capped by the same band slope: the constancy of \(c\) is a
\emph{graph property}, not a postulate. \textbf{{[}measured{]}} Vacuum
optics at the adopted law table are linear (intensity-independent) and
dispersive exactly per this band.

\section{The door pair: how field becomes
matter}\label{the-door-pair-how-field-becomes-matter}

\subsection{Why conversion needs a
register}\label{why-conversion-needs-a-register}

The door of §2.4 says conversion is quantized. But \emph{refraction} ---
light slowing in glass --- requires the field in matter to be
continuously, reversibly, \emph{partially} converting into a matter-side
response and back (this is the actual textbook mechanism of the
refractive index: the field polarizes the medium, the polarization
re-radiates). A conversion that is only ever whole atoms cannot do this.
So the substrate needs a state for \emph{incomplete conversion to stand
in}: the register \(\chi\). Energy rotates continuously and reversibly
between \(\psi_e\) and \(\chi\) (``virtual'' conversion --- optics), and
settles from \(\chi\) into matter only in whole atoms at discrete clicks
(``real'' conversion). §8 proves this two-stage structure is not a
convenience but a necessity.

\subsection{The pair}\label{the-pair}

Between clicks, the field and the register in one cell obey the coupled
first-order pair \textbf{{[}derived; implemented in v94; conservation of
\(E_e + |\chi|^2\) under the pair is exact{]}} \[
i\,\frac{d\psi_e}{dt} = \Omega_1\,\psi_e - g\,\chi, \qquad
i\,\frac{d\chi}{dt}   = \Omega_2\,\chi   - g\,\psi_e,
\] with \(\Omega_2 = w_{2e}\) (the register is matter-pitched: it is
matter in the making) and coupling \(g = \mathrm{door}_g\, w_2\), built
from the bare matter pitch --- the door's stiffness belongs to the
fabric, not the local load. Note the shape: \emph{this is the Maxwell
pair again}, pointing along the mode direction instead of the space
direction. Intake (field \(\to\) register) and outtake (register \(\to\)
field) are the two half-cycles of one rotation; conversion is incomplete
by construction, with maximum transferred fraction
\[F(x) = \frac{4g^2}{\Delta(x)^2 + 4g^2}, \qquad
  \Delta(x) = \Omega_1 - \Omega_2 = \frac{w_1 - w_2}{1 + q\,x}.\]
\(\Delta\) is the \textbf{detuning} --- the mismatch between the two
clocks. In vacuum (\(x=0\)) the mismatch is maximal and the door barely
opens; in loaded matter the flattening law shrinks \(\Delta\) and the
door opens. \emph{``Empty space demands a far threshold; bound matter
converts freely''} --- the work-function rule --- stops being a rule and
becomes clock arithmetic. \textbf{{[}derived from the adopted pitch
law{]}}

\subsection{The beat, the winding, and the
click}\label{the-beat-the-winding-and-the-click}

\begin{figure}[t]
\centering
\includegraphics[width=\linewidth]{fig_click.pdf}
\caption{\textbf{The click grain.} Top: the accumulated phase difference
between the field and matter clocks (the beat). The door may fire only
when it completes a full turn ($2\pi$) --- a \emph{winding} event.
Bottom: between clicks the register $|\chi|^2$ evolves continuously and
reversibly; matter $E_m$ changes only at clicks, and only by whole atoms
$\varepsilon$. Illustration with law-table values.}
\label{fig:click}
\end{figure}

\emph{When} may the door fire? The kernel keeps, per cell, the
accumulated phase difference between the two clocks:
\[\varphi_D \mathrel{+}= (w_{1e} - w_{2e})\,dt,\] the \textbf{beat}. The
door fires only when \(\varphi_D\) completes a full turn
(\(|\varphi_D| = 2\pi\)) --- and then transfers whole atoms
\(\varepsilon\) from the register into matter, by an energy-preserving
rotation between \(\chi\) and \(\psi_m\) solved to move exactly the
fired amount (Fig. \ref{fig:click}). Three facts about this arrangement,
each doing real work:

\begin{enumerate}
\def\labelenumi{\arabic{enumi}.}
\tightlist
\item
  \textbf{The click is a topological event.} A full turn of a phase is a
  \emph{winding}, and winding numbers are integers by geometry, not by
  decree. The integer character of conversion --- you get 0, 1, 2 atoms,
  never \(1.37\) --- comes for free from topology. Only the \emph{size}
  of the atom, \(A_0\), is a postulate (§8).
\item
  \textbf{The click clock is load-throttled.} The beat rate is
  \(|\Delta(x)| = (w_2 - w_1)/(1+qx)\): loaded cells click \emph{less
  often}. This derived throttle is one layer of the substrate's defence
  against runaway conversion (§5.3).
\item
  \textbf{Phase can cross the door.} The transfer is a rotation, not an
  erasure: the register's clock reading composes into matter's clock
  algebraically. Whether it \emph{survives} the environment is a
  measured question --- and the current answer on one geometry is
  \emph{no} (§10, the retention wall).
\end{enumerate}

\subsection{The mass--energy
correspondence}\label{the-massenergy-correspondence}

\begin{quote}
\textbf{The familiar law.} \(E = mc^2\): mass is a form of energy; a
mass \(m\) at rest ``contains'' energy \(mc^2\), and nuclear reactions
that destroy mass release exactly that energy.
\end{quote}

In the substrate this famous equation is not a surprising discovery but
a \emph{dissolved identity}. Matter simply \textbf{is} energy in the
dense mode: \(E_m\) is one ledger column. ``Mass'' is the name for how
much energy sits settled in matter; ``creating matter'' is moving ledger
entries through the door; ``annihilation'' is the door running the other
way. In substrate units where \(c_{\mathrm{sub}} = 1\), the relation
\(E = m\) is the \emph{definition} of the ledger, and destroying a mass
without its energy appearing elsewhere is not merely forbidden --- it is
not even expressible.

\textbf{A worked ledger line, with law-table numbers.} Take one loaded
cell and let the door fire one atom of matter outward. The matter column
drops by exactly
\(\varepsilon = A_0\,w_{2e}/2\pi \approx 1.15 \times 2.09 / 6.28 \approx 0.38\)
energy units; the field column rises by the same \(0.38\); the
conservation meter does not move. That \emph{is} the content of ``mass
converted to energy'' here: one number moving between two columns of a
ledger that always sums to the same total. And the \(c^2\)? In substrate
units the wave speed is \(1\) and \(E = m\) outright; the \(c^2\) of the
SI formula is a \emph{unit conversion} --- the price of measuring mass
in kilograms but energy in joules --- not a physical ingredient.

\textbf{What \(E = mc^2\) \emph{does} give us here, measured:}

\begin{itemize}
\tightlist
\item
  conversion is quantized (the atoms) --- \textbf{{[}measured{]}};
\item
  energy leaving matter arrives in the field at the right total ---
  \textbf{{[}measured, drift floor{]}}.
\end{itemize}

\textbf{What it does \emph{not} yet give us:}

\begin{itemize}
\tightlist
\item
  the \emph{inertial} content of mass --- that momentum crosses the door
  with the energy --- is partially \textbf{{[}open{]}}: momentum is the
  first moment of conversion (a centre-of-energy theorem replaces
  Noether's theorem on the roster), but the measured radiation-pressure
  delivery at the conversion door is roughly \(100\times\) below the
  naive expectation. The v94 door pair predicts this deficit shrinks
  because phase (and hence the momentum-carrying gradient) now crosses
  the door as algebra; that is a registered, scored prediction, not yet
  a result.
\end{itemize}

\section{The metabolic cycle: matter that eats and
shines}\label{the-metabolic-cycle-matter-that-eats-and-shines}

\subsection{Intake, outtake, and the
polariton}\label{intake-outtake-and-the-polariton}

\begin{figure}[t]
\centering
\includegraphics[width=\linewidth]{fig_dispersion.pdf}
\caption{\textbf{The mixed branch.} Coupling the travelling field band
(blue dashed) to the stationary register line (grey dashed) splits the
spectrum into two mixed branches $\omega_\pm$ separated by $2g$ at the
crossing. A wave on a mixed branch is part field, part register: it
travels at $c_{\mathrm{sub}}$ times its field fraction --- slower in
matter. This is refraction, derived from the door.}
\label{fig:dispersion}
\end{figure}

With the pair in place, an object's life is a cycle: \textbf{intake}
(ambient field settles through the door into \(E_m\)) and
\textbf{outtake} (matter returns energy fieldward --- spontaneously as
radiance, §5.2, or stimulated, through the same pair run backward).
Coupling the travelling field to the stationary register splits the
spectrum into \textbf{two mixed travelling waves} (Fig.
\ref{fig:dispersion}; physicists call such a mixed light--matter wave a
\emph{polariton}; the branch formula is in Appendix C).
\textbf{{[}derived{]}} Away from the crossing one branch is field-like
(speed \(c_{\mathrm{sub}}\)) and one register-like (going nowhere); near
it, a wave is \emph{part field, part matter-in-the-making}, and it
travels at \[v = c_{\mathrm{sub}} \cdot W_e,\] where \(W_e < 1\) is
simply \emph{the fraction of the mixed wave that is still field}. A wave
that spends part of each cycle as register does not spend that part
moving --- the refractive slow-down, now a consequence of the door
rather than a missing mechanism. The same algebra predicts a
\textbf{metabolic wave}: a disturbance of the intake/outtake cycle
propagating at a finite, \(g\)-tunable speed \(v_M\) (the in-flight
experiment M1 measures the \(v_M(\mathrm{door}_g)\) curve).

\subsection{\texorpdfstring{Radiance and the \(T^4\)
correspondence}{Radiance and the T\^{}4 correspondence}}\label{radiance-and-the-t4-correspondence}

\begin{quote}
\textbf{The familiar law.} A hot body radiates with power
\(P = \sigma A T^4\) (Stefan--Boltzmann): double the absolute
temperature, sixteen times the light.
\end{quote}

\begin{figure}[t]
\centering
\includegraphics[width=\linewidth]{fig_radiance.pdf}
\caption{\textbf{The metabolic balance.} Outtake is graded steeply in the
load ($O \propto x^4$, orange); intake opens with load but is
supply-limited (blue). Where they cross with outtake steeper, a
\emph{restoring fixed point} $x^*$ exists: matter that eats and shines
and holds its size. The park itself is the in-flight measurement M6.}
\label{fig:radiance}
\end{figure}

Before the formula, its status, so no one hunts for a missing
derivation: the \emph{form} below --- a steep, graded outtake --- is
\textbf{adopted}: chosen, then held fixed under the battery. Whether the
substrate \emph{must} use the fourth power is open; the mode-counting
sketch afterwards is the candidate derivation, no more. With that said:
every beat, a loaded cell sheds
\[O(x) = k_{\mathrm{rad}}\;\mathrm{cap}\;x^{\,p_{\mathrm{rad}}},
  \qquad p_{\mathrm{rad}} = 4,\] into the field, in atoms, where
\(x = E_m/\mathrm{cap}\) is the holdings load
(\(k_{\mathrm{rad}} = 0.05\) in the law table). The correspondence: the
load \(x\) --- how full, hence how fast-clocked and crowded, a cell is
--- is the substrate's temperature-\emph{like} variable (a fullness, not
a thermometer reading --- the analogy is structural), and the quartic
grading is the \(T^4\) shape. \emph{Why} the fourth power? The
mode-counting sketch: in equilibrium radiation in \(D\) spatial
dimensions, energy density scales as \(T^{D+1}\) (each of the
\(\sim T^D\) available field modes carries \(\sim T\)); on the
three-dimensional contact graph this gives exponent \(4\)
\textbf{{[}open{]}}. What the law demonstrably buys
\textbf{{[}measured{]}}: hoarded, non-metabolizing energy blobs are
taxed away (half-lives drop from \(\sim 260\)--\(510\) to
\(80\)--\(140\) time units when radiance is on), a bath of cells
acquires a glow, and the balance diagram of Fig. \ref{fig:radiance}
acquires its restoring wall.

\subsection{The measured intake ramp, the fold, and why nothing
ignites}\label{the-measured-intake-ramp-the-fold-and-why-nothing-ignites}

\begin{figure*}[t]
\centering
\includegraphics[width=0.92\textwidth]{fig_fold.pdf}
\caption{\textbf{Virtual dressing has a threshold structure.} (a)
Predicted: the self-consistent register dressing $u$ versus drive $\beta$
has a fold --- beyond it no small reversible solution exists.
(b) Measured (M3b ramp, v94): door intake versus per-cell field energy.
Intake wakes exactly at one atom ($E_e \approx 0.36 \approx \varepsilon$),
the coherent channel dominates from $E_e \approx 1.0$ (the fold, within
its $O(1)$ calibration), stimulated outtake wakes above that and, with
capacity defence, keeps the growth bounded: no runaway to the highest
drive tested.}
\label{fig:fold}
\end{figure*}

A door that opens with load invites a worry: coherent light shining on
matter makes it grow; growth opens the door further; does matter ignite
and eat the world? The substrate defends in layers, all structural
\textbf{{[}derived{]}}, all now probed \textbf{{[}measured{]}}:

\begin{enumerate}
\def\labelenumi{\arabic{enumi}.}
\tightlist
\item
  vacuum is maximally detuned (the door barely opens where there is
  nothing);
\item
  a random-phase bath drives the register incoherently --- its pushes
  average out (no net pumping without phase order);
\item
  real conversion is gated: below one atom, nothing fires (Fig.
  \ref{fig:fold}b: intake wakes \emph{exactly} at
  \(E_e \approx \varepsilon\)); a presence rule makes truly empty cells
  demand an unpayable threshold (vacuum cannot settle at any tested
  intensity);
\item
  the pair is a rotation --- it can redistribute energy but never
  amplify it;
\item
  loaded cells click more slowly (the beat throttle, §4.3).
\end{enumerate}

Above a derived drive threshold --- the \textbf{fold} of Fig.
\ref{fig:fold}a, meaning the point beyond which no \emph{small,
reversible} response exists; predicted from the pair plus the flattening
law, and located within its stated \(O(1)\) factor by measurement ---
the register dressing over \emph{existing} matter outgrows the
small-response regime and conversion accelerates (\(\propto E_e^{1.75}\)
measured). Even there: \textbf{no runaway} up to the strongest drive
tested --- capacity defence (over-full cells evaporate) and the
stimulated outtake bound the growth, with the energy ledger at the
round-off floor throughout. The fold is thus not an instability but the
substrate's \emph{stimulated conversion threshold} --- the intensity at
which light over matter breeds matter, boundedly: pair production's
structural cousin. \textbf{{[}measured, v94 M3b{]}}

\section{Gravity: space flow}\label{gravity-space-flow}

\subsection{Matter digs a well in the space
ledger}\label{matter-digs-a-well-in-the-space-ledger}

\begin{figure}[t]
\centering
\includegraphics[width=\linewidth]{fig_gravity.pdf}
\caption{\textbf{The footprint.} A mass is converted space: its cells run
a maintained deficit of $E_s$, graded outward over a few cells
\textbf{[measured, g1]}. Space transport is pressure-only --- flow runs
down the gradient toward the deficit; nothing attracts at a distance. The
steady $1/r$ far field of Newtonian gravity has \emph{not} been measured
(g4) and is an open problem tied to the metabolic cycle.}
\label{fig:gravity}
\end{figure}

Recall the founding sentence: matter is \emph{converted space}. Where
matter condensed, space-energy was consumed. A live mass therefore sits
in a maintained, graded depression of \(E_s\) --- its \textbf{footprint}
--- extending over its neighbourhood (Fig. \ref{fig:gravity})
\textbf{{[}measured, experiment g1{]}}. Space transport in the substrate
obeys one signed rule: \textbf{pressure pushes; nothing pulls}. Cells
with more \(E_s\) push space toward cells with less; ``attraction'' is
only ever the appearance of being pushed down someone's gradient.
Consequences, each a gated measurement:

\begin{itemize}
\tightlist
\item
  \textbf{Opacity (g3):} matter is emergently opaque to space flow ---
  occultation happens; a mass shadows what is behind it.
  \textbf{{[}measured{]}}
\item
  \textbf{The flat-medium result} (internally,
  ``\(\pi\)-flatness'')\textbf{:} matter with \emph{closed books}
  (intake = outtake, nothing net flowing) loads the surrounding medium
  with essentially nothing: measured medium response flat to a few parts
  in \(10^{4}\) across a ladder of stable objects. Only \emph{dying}
  (net-radiating) matter disturbs the medium at range.
  \textbf{{[}measured, v91{]}}
\end{itemize}

\subsection{The Poisson correspondence --- and its honest
limit}\label{the-poisson-correspondence-and-its-honest-limit}

\begin{quote}
\textbf{The familiar law.} Newtonian gravity, in field form: the
potential \(\Phi\) obeys Poisson's equation
\(\nabla^2\Phi = 4\pi G\rho\), giving the \(1/r\) potential of a point
mass, and bodies accelerate down \(-\nabla\Phi\).
\end{quote}

The near-field correspondence is direct and needs no jargon: \emph{a
permanent drain makes a permanent dip.} Matter converting space is the
drain; the maintained \(E_s\) deficit is the dip, and its shape is set
by how space pushes downhill toward it. In the map to Newton: the depth
of the dip plays the potential \(\Phi\); the conversion rate plays the
mass density \(\rho\); being pushed down someone's dip plays
gravitational attraction. Down-gradient drift of nearby objects is
measured. \textbf{{[}measured at footprint range{]}}

The far field is where honesty bites. Newtonian gravity requires the
\(1/r\) tail to extend \emph{arbitrarily far, statically}. The substrate
measurement (experiment g4) is negative so far: a leaking blob's entire
space flux is accounted by its mass-rate bookkeeping --- outward wind
balanced by inward refill --- with \textbf{no steady \(1/r\) tail} (a
``monopole'' tail, in the jargon) \textbf{{[}measured negative{]}}. The
theory's own diagnosis: a \emph{static} far field must be maintained by
a \emph{cycle} --- a stable particle whose internal metabolism
continuously turns space over --- and no such particle has yet parked
(the M6 fixed point is in flight). Newton's far field is therefore
\textbf{{[}open{]}}, and marked as such in the correspondence table: the
substrate currently derives the well, the shadow, and the flatness
theorem, but not yet the \(1/r\) tail. A theory that quietly
extrapolated the footprint to infinity would be lying with a graph.

\section{Interference: the double slit, click by
click}\label{interference-the-double-slit-click-by-click}

\begin{quote}
\textbf{The familiar experiment.} Light (or single electrons!) sent
through two narrow slits lands on a far screen in \emph{fringes} ---
bands of bright and dark. The pattern persists even when quanta arrive
one at a time: each click is pointlike, yet thousands of clicks trace
the wave pattern. This is the standard demonstration that nature adds
\emph{amplitudes}, not probabilities.
\end{quote}

\begin{figure}[t]
\centering
\includegraphics[width=\linewidth]{fig_doubleslit.pdf}
\caption{\textbf{Fringes, twice.} Top: the wave intensity on the far
screen --- the two path families' arrows added, then squared. Bottom: 362
single-atom clicks sampled by the door (the count matching the v89/v90
tier-1 run), rebuilding the same fringes one click at a time. Measured
minima on the substrate are $7.7\times$ darker than neighbouring maxima,
at parameter-free predicted positions.}
\label{fig:doubleslit}
\end{figure}

The substrate derivation needs nothing beyond §3 and §4:

\begin{enumerate}
\def\labelenumi{\arabic{enumi}.}
\tightlist
\item
  \textbf{Transport carries phase.} The field update is \emph{unitary}
  --- it preserves arrow lengths and relative clock angles; a packet
  splitting through two openings arrives at a far cell as the \emph{sum
  of two arrows}, one per path family, each having precessed through its
  own path length: phase difference \(\Delta\varphi = k\,\Delta L\).
\item
  \textbf{Energy is the squared sum.} The far cell's field energy is
  \(|\psi_1 + \psi_2|^2 = |\psi_1|^2 + |\psi_2|^2 + 2|\psi_1||\psi_2|\cos\Delta\varphi\)
  --- the interference term is just the cross-term of adding arrows
  before squaring. Maxima where paths agree, near-zeros where they
  cancel (Fig. \ref{fig:doubleslit}, top).
\item
  \textbf{Detection is the door.} A detector cell converts field to
  matter through the same atom-grained door as everything else: whole
  clicks, at a rate set by the local field energy. Sparse arrivals
  therefore sample the fringe pattern probabilistically --- detection
  probability tracks \(|\psi|^2\). (In quantum mechanics that rule
  carries the name \emph{Born rule}; here it is door behaviour, not an
  extra axiom.)
\end{enumerate}

The measured record \textbf{{[}v89--v90, battery-gated{]}}: fringes at
parameter-free predicted loci; detection statistics matching
\(|\psi|^2\) to \(\pm 0.02\); and the tier-1 claim --- \textbf{362
whole-atom clicks, arriving one at a time, rebuilding the fringes} with
minima \(7.7\times\) dark and the correct fringe phase (coherence score
0.683 versus 0.380 for the control). Eraser fringes/anti-fringes,
delayed choice, and exact no-signaling are further gated tiers of the
same apparatus. The double slit is not an interpretation puzzle here; it
is what an atom-doored wave substrate \emph{does}.

\section{The unification theorem}\label{the-unification-theorem}

\subsection{One primitive, three edge
types}\label{one-primitive-three-edge-types}

Follow the arc backwards. Spatial field links (the \emph{field hop})
move field by precession + exchange rotation (§3). Spatial matter links
(the \emph{matter hop}) move matter by the \emph{same} algebra. The v94
door pair showed the \emph{conversion door} is a third link type --- a
link in the \textbf{mode} direction rather than the space direction ---
carrying the same primitive, with one distinguishing feature: its
settlement is quantized. One graph, one rotation algebra, three edge
families (Fig. \ref{fig:substrate}b). That is the v94 unification of
\emph{law}.

\subsection{Can matter be retired
altogether?}\label{can-matter-be-retired-altogether}

The obvious next question --- the deepest the programme has asked --- is
whether the unification can also erase the \emph{structure}: could
matter be just a self-bound lump of the one field, with \(\psi_m\) (and
the click) eliminated? The answer is \textbf{no, by a theorem} --- and
the theorem is more informative than the hope.

In plain language first: \emph{if light really slows in matter by
continuously dipping partway into conversion and back (that is what
refraction is), and if settled matter only ever changes by whole packets
(that is what the measurements say), and if the rules are local with no
hidden eternal grid, then every cell needs} \textbf{two}
\emph{conversion stages with a temporary holding slot between them. One
stage alone cannot be both smooth enough for glass and chunky enough for
clicks.} The formal version (the ``fiber'' is just the cell's internal
conversion path, field--register--matter):

\textbf{Theorem (fiber minimality).} \emph{Assume:}

\begin{itemize}
\tightlist
\item
  \textbf{(H1) Refraction is real:} loaded matter slows the field by
  continuous, reversible, phase-coherent partial conversion (§4.1, §5.1
  --- the mechanism behind every refractive index);
\item
  \textbf{(H2) The grain is real:} settled matter changes only in whole
  atoms \(\varepsilon = A_0\omega/2\pi\) (measured to \(\sim 10^{-8}\));
\item
  \textbf{(H3) Dynamics is local and memoryless:} the update is per-cell
  / per-link, one step at a time, with no persistent hidden reference
  (no background --- constraint 2).
\end{itemize}

\emph{Then} the conversion path of every cell must contain \textbf{at
least two connections} --- one continuous, carrying \emph{no}
settled-inventory change, and one firing only at discrete instants,
carrying \emph{all} of it, in atoms --- \textbf{and a persistent
two-component carrier between them, distinct from the settled
inventory.} That is: \(\psi_e\,\)---\(\,\chi\,\)---\(\,\psi_m\), exactly
the three-node fiber of Fig. \ref{fig:substrate}a, is minimal.

\emph{Proof sketch.} A single connection would have to transfer energy
both continuously (H1: the dressing exists between clicks) and in exact
integers into the inventory (H2). Continuous transfer with integer
inventory steps means the connection's two ends do not form a conserved
pair between clicks --- contradicting the rotation primitive --- unless
the transfer is zero between clicks, contradicting H1. Hence two
connections with the stated roles. The node between them must hold
in-transit energy \emph{and its clock reading} across many steps (the
reversible dressing persists over many beats), hence a persistent
two-component amplitude; identify it with the inventory and H2 applies
to it, killing H1 again. \(\square\)

The escapes are exactly: drop H2 (contradicts measurement), drop H3
(hidden memory --- a background by another name), or quantize the
substrate itself (a different theory). The register is not scaffolding;
it is forced.

\subsection{What remains imposed}\label{what-remains-imposed}

The theorem sharpens what the door's ``specialness'' actually consists
of. Three irreducibles:

\begin{enumerate}
\def\labelenumi{\arabic{enumi}.}
\tightlist
\item
  the \textbf{two-connection fiber} --- imposed, and now \emph{proven
  minimal};
\item
  the click's \textbf{timing and integer character} --- \emph{not}
  imposed: derived from winding topology (§4.3);
\item
  the \textbf{action scale \(A_0\)} --- imposed. This single constant is
  the substrate's Planck postulate: the one quantum assumption the
  theory makes. Deriving \(A_0\) would require quantizing the
  substrate's own dynamics --- a different programme, deliberately not
  this one.
\end{enumerate}

\section{The correspondence table}\label{the-correspondence-table}

\begin{table*}[t]
\centering
\small
\begin{tabular}{p{2.6cm} p{5.4cm} p{2.2cm} p{5.6cm}}
\hline
\textbf{Familiar law} & \textbf{SCP mechanism} & \textbf{Status} &
\textbf{Measurement (v89--v94)} \\
\hline
Maxwell pair; fixed $c$ &
two first-order mutually-inducing quadratures; $c_{\mathrm{sub}}$ = band
slope of the roster graph (\S3) &
exact in shape; scalar (no vector curl) &
linear vacuum optics; dispersion per band; conservation at
$10^{-13}$--$10^{-15}$ \\
Refraction ($v<c$ in matter) &
polariton mixing of field with the door register (\S5.1) &
derived (v94); dressing measured &
register filling and mixed-channel intake measured on the M3b ramp;
$v_M$ curve in flight (M1) \\
$E=mc^2$ &
matter \emph{is} the dense-mode ledger column; conversion through the
atom-grained door (\S4.4) &
identity (ledger); inertia part open &
atom linearity $\sim 10^{-8}$; energy books close at drift floor;
radiation-pressure deficit $\sim 100\times$ recorded, engaged by the
door pair (scored prediction) \\
Stefan--Boltzmann $T^4$ &
graded outtake $O = k_{\mathrm{rad}}\,\mathrm{cap}\,x^4$ per beat; $x$
the temperature-like load (\S5.2) &
form adopted + gated; exponent derivation open &
hoard taxing (half-lives $260$--$510 \to 80$--$140$); bath glow; battery
87/87 at the adopted table \\
Poisson / Newton &
maintained $E_s$ footprint = Green's-function well; pressure-push
transport (\S6) &
near field derived + measured; far field OPEN &
g1 footprint; g3 opacity; $\pi$-flatness $\sim 10^{-4}$; g4: no steady
$1/r$ monopole (negative) \\
Double slit / Born rule &
unitary phase transport; arrows added then squared; detection = the atom
door (\S7) &
derived; measured &
fringes at parameter-free loci; Born to $\pm0.02$; 362 single clicks
rebuild fringes, minima $7.7\times$ dark; no-signaling exact \\
Quantization ($\hbar$) &
click timing + integers from beat-phase winding topology; scale $A_0$
postulated (\S4.3, \S8) &
timing/integers derived; $A_0$ imposed &
$\varepsilon = A_0\omega/2\pi$ linearity $\sim 10^{-8}$; measured intake
onset exactly at one atom (M3b) \\
\hline
\end{tabular}
\caption{Where the familiar laws come from on the substrate, with the
grading of each claim and the experiment that gates it.}
\label{tab:corr}
\end{table*}

Table \ref{tab:corr} collects the derivations of §3--§7 with their
status and their gating measurements. The pattern worth noticing: the
correspondences are strongest exactly where the substrate's one
primitive operates unobstructed (waves, interference, conversion
bookkeeping), and open exactly where a \emph{persistent structure} must
hold itself together against its environment (the far field needs a
parked particle; inertia needs phase across the door; see §10). The
theory's frontier is not a missing force --- it is the stability of
matter.

\section{Open questions}\label{open-questions}

The programme's rule is that the open list is load-bearing. In current
priority order:

\begin{enumerate}
\def\labelenumi{\arabic{enumi}.}
\tightlist
\item
  \textbf{The retention wall (M4, measured negative).} On a test ring
  carrying a quantized circulation, door traffic \emph{erodes} the phase
  structure it should carry: the stimulated outtake samples an
  oscillating cross-term at one frozen instant per beat, and cells
  frozen on the draining half fire every beat --- a ratchet. The
  diagnosis is precise (the sampling, not the pair, is at fault) and the
  registered fix is principled: integrate the cross-term over the full
  beat (synchronous detection), so reversible oscillation averages to
  zero and only genuinely phase-locked geometry drives outtake. Notably,
  the same measurement produced a positive surprise: on a hot bath the
  ring \emph{self-heals} its circulation without the door pair ---
  topology plus accretion is itself a phase-retention mechanism.
\item
  \textbf{The metabolic fixed point (M6, in flight).} Does matter park
  --- intake = outtake at a restoring load \(x^*\) (Fig.
  \ref{fig:radiance})? This is the programme's oldest open ledger and
  the prerequisite for the g4 far field.
\item
  \textbf{The metabolic wave (M1) and the mobility release (M5), in
  flight.} Does the cycle propagate at the predicted \(v_M\), and can an
  object bound by self-trapping \emph{move by converting} --- intake at
  the front, outtake at the rear --- when its internal transport is
  locked?
\item
  \textbf{The far field (g4).} No steady \(1/r\) tail without a parked,
  metabolizing particle; depends on 2.
\item
  \textbf{The action scale \(A_0\).} The one imposed quantum constant
  (§8.3). Whether any deeper substrate could yield it is outside the
  current programme's rules.
\item
  \textbf{The path to atoms.} Shells, exclusion, chemistry: nothing
  beyond suggestive structure yet; blocked behind stable particles (2)
  and phase retention (1).
\item
  \textbf{Engineering debt, stated because honesty includes plumbing:}
  the two independent kernel implementations agree bit-for-bit on the
  gated battery but diverge at strength on one transport path; all
  campaign data therefore come from the designated campaign kernel until
  the debt is paid.
\end{enumerate}

\section{Appendix A: notation}\label{appendix-a-notation}

Table \ref{tab:notation} collects every symbol used in the paper, with
the adopted law-table value where one exists.

\begin{table}[h]
\centering
\small
\begin{tabular}{p{1.5cm} p{6.0cm}}
\hline
\textbf{Symbol} & \textbf{Meaning} \\
\hline
$\psi_e$ & field amplitude per cell ($=f_{a1}+i f_{a2}$); energy $E_e = |\psi_e|^2$ \\
$\psi_m$ & matter amplitude; energy $E_m = |\psi_m|^2$, clock $\theta_2$ \\
$\chi$ & door register: energy in mode-transit; $|\chi|^2$ booked in the ledger \\
$E_s$ & space energy (sets cell size); non-oscillating \\
flight & energy in space-transit on a link (mid-hop parcels) \\
$w_1, w_2$ & bare pitches of field and matter ($1.65$, $2.9$) \\
$x$ & load: (holdings + transit)/$\mathrm{cap}$ \\
$q$ & pitch-flattening constant ($1.2$); $w_{e} = w/(1+qx)$ \\
$\mathrm{cap}$ & cell carrying capacity ($2.5$) \\
$\Omega_1, \Omega_2$ & effective pitches $w_{1e}, w_{2e}$ \\
$\Delta(x)$ & detuning $\Omega_1 - \Omega_2$; the beat rate \\
$\varphi_D$ & accumulated beat phase; door fires at $|\varphi_D| = 2\pi$ \\
$g$ & door coupling $= \mathrm{door}_g \cdot w_2$ (knob $\mathrm{door}_g$) \\
$F(x)$ & maximum transfer fraction $4g^2/(\Delta^2 + 4g^2)$ \\
$A_0$ & action atom (the substrate's $\hbar$; $1.15$) \\
$\varepsilon$ & energy atom $A_0\omega/2\pi$ at source pitch $\omega$ \\
$K_{ij}$ & link exchange rate; $\bar K$ its mean \\
$\epsilon(k)$ & field band $2\bar K(1 - \cos kd)$ \\
$c_{\mathrm{sub}}$ & substrate light speed $\approx 2\bar K d$ \\
$k_{\mathrm{rad}}, p_{\mathrm{rad}}$ & radiance constants ($0.05$, $4$): $O = k_{\mathrm{rad}}\,\mathrm{cap}\,x^{p_{\mathrm{rad}}}$ \\
$\omega_\pm(k)$ & polariton branches of the door pair \\
$v_M$ & metabolic wave speed (mixed-branch group velocity) \\
\hline
\end{tabular}
\caption{Notation. All quantities in substrate units (cell spacing, step
time); no SI constants are imported anywhere in the theory.}
\label{tab:notation}
\end{table}

\section{Appendix B: method --- the battery and the
ratchet}\label{appendix-b-method-the-battery-and-the-ratchet}

Every claim tagged \textbf{{[}measured{]}} is gated by a regression
battery: a fixed suite of experiments (87 at the current law table;
grown from 20 at v89) that every modification of the substrate must
re-pass \emph{in full} before adoption. Bars encode claims
quantitatively and may be sharpened by measurement but never softened to
pass; a green test leaves the suite only by explicit decision. New
mechanisms enter behind knobs that are \textbf{byte-inert at default
zero} --- with the knob off, the simulation reproduces the previous law
table bit-for-bit --- so that every comparison is against a frozen,
reproducible baseline. Two independent kernel implementations (C and Go)
mirror the update; the gated battery is bit-identical between them.
Conservation is not asserted but \emph{metered}: the ledger of Appendix
A is summed every step, and its drift is a tripwire (it caught a real
bookkeeping bug during v94 face-2 testing at eight orders of magnitude
above the floor). The theory in this paper is, in that specific sense,
an \emph{instrumented} theory: each section's claims name the experiment
that would have falsified them.

\section{Appendix C: the mixed-branch
formula}\label{appendix-c-the-mixed-branch-formula}

The two mixed travelling waves of §5.1 (the avoided crossing of Fig.
\ref{fig:dispersion}) have frequencies, as functions of wavenumber
\(k\):

\[
\omega_\pm(k) = \tfrac{1}{2}\big[\Omega_1 + \epsilon(k) + \Omega_2\big]
 \pm \sqrt{\tfrac{1}{4}\big[\Omega_1 + \epsilon(k) - \Omega_2\big]^2 + g^2}.
\]

This is the standard result of coupling a travelling band
\(\Omega_1 + \epsilon(k)\) to a stationary line \(\Omega_2\) with
strength \(g\): far from the crossing each branch follows its parent; at
the crossing they repel by \(2g\) and each is an equal mix. The field
weight \(W_e(k)\) of §5.1 is the squared field share of the branch's
makeup, and the branch's speed is the parent band's speed times that
share.
